
\setcounter{subsection}{6}% sincronización del rótulo visible XIII.4.6.0 de la fuente
\setcounter{proposition}{5}% sincronización con la numeración visible XIII.4.6 de la fuente
\begin{proposition}[Fórmula de Künneth]
\label{XIII.4.6}
\index{Kunneth (formula de)@K\"unneth (fórmula de)|hyperpage}%
Sean $k$ un cuerpo separablemente cerrado de característica
$p\geq 0$, $X$ e~$Y$ dos $k$-esquemas conexos, $a$ un punto
geométrico de~$X$, $b$ un punto geométrico de~$Y$ y~$c$ un punto
geométrico de~$X\times_k Y$ situado sobre~$a$ y~$b$. Se supone que se
satisface una de las dos condiciones siguientes:
\begin{enumerate}
\item[a)] $X$ es de tipo finito sobre~$k$ y los esquemas de tipo finito
sobre una clausura algebraica~$\bar{k}$ de~$k$, de dimensión
$\leq\dim X$, son fuertemente desingularizables. \textup{(SGA~5 I~3.1.5)}.
\item[b)] $X$ es cuasicompacto y cuasiseparado y todo esquema de tipo
finito sobre~$\bar{k}$ es fuertemente desingularizable.
\end{enumerate}

Entonces, si $p'$ es el conjunto de los números primos distintos de~$p$,
el morfismo
\begin{equation*}
\label{eq:XIII.4.6.0} \tag{\thesubsection.0} {}
\pi_1^{p'}(X\times_k Y,c)\to\pi_1^{p'}(X,a)\times\pi_1^{p'}(Y,b),
\end{equation*}
deducido
\marginpar{422}
de los homomorfismos de grupos fundamentales asociados a las proyecciones
$$
X\times_k Y\to X \quad \text{y} \quad X\times_k Y\to Y,
$$
es un isomorfismo.
\end{proposition}

Puede suponerse que~$k$ es algebraicamente cerrado y que~$X$ es reducido
(SGA~4 VIII~1.1). Sean $Z=X\times_k Y$, $g\colon X\to k$ y $f\colon
Z\to Y$ los morfismos canónicos. El morfismo~$g$, y por tanto también~$f$,
es universalmente localmente $1$-asférico para~$p'$
según~
\ifthenelse{\boolean{orig}}
{\Ref{XIII.3.5}}
{\Ref{XIII.3.4}}.
Como~$X$ es conexo, $f$ es
$0$-acíclico (SGA~4 XV~1.16). Por otra parte, de~\Ref{XIII.3.2} resulta
que, para todo $p'$-grupo finito~$C$, $(C_X,g)$ es cohomológicamente propio
relativamente a~$k$ en dimensión $\leq 1$. De ello resulta que~$(C_Z,f)$
es cohomológicamente propio relativamente a~$Y$ en dimensión $\leq 1$
(\Ref{XIII.1.5}~c)) y que $f_*C_Z$ y $\R^1f_*C_Z$ son haces constantes.
Por consiguiente,
\ifthenelse{\boolean{orig}}
{$g$}
{$f$}
satisface todas las hipótesis de~\Ref{XIII.4.2}. Se tiene, pues, la sucesión
exacta
$$
\pi_1^{p'}(X_b,c)\to\pi_1^{p'}(Z,c)\to\pi_1^{p'}(Y,b)\to 1.
$$
Además, el morfismo compuesto
$$
\pi_1^{p'}(X_b,c)\to\pi_1^{p'}(Z,c)\to
\ifthenelse{\boolean{orig}}
{\pi_1^{p'}(Y,b)}
{\pi_1^{p'}(X,a)}
$$
es un isomorfismo y se tiene, por tanto, la sucesión exacta
$$
1\to\pi_1^{p'}(X,a)\to\pi_1^{p'}(Z,c)\to\pi_1^{p'}(Y,b)\to 1.
$$
Por otra parte, el morfismo~\eqref{eq:XIII.4.6.0} permite definir
un morfismo de esta sucesión exacta en la sucesión exacta
$$
1\to\pi_1^{p'}(X,a)\to
\pi_1^{p'}(X,a)\times\pi_1^{p'}(Y,b)
\ifthenelse{\boolean{orig}}
{\times}
{\to}
\pi_1^{p'}(Y,b)\to 1,
$$
y de ello resulta que el morfismo~\eqref{eq:XIII.4.6.0} es un
isomorfismo.
